%% Согласно ГОСТ Р 7.0.11-2011:
%% 5.3.3 В заключении диссертации излагают итоги выполненного исследования, рекомендации, перспективы дальнейшей разработки темы.
%% 9.2.3 В заключении автореферата диссертации излагают итоги данного исследования, рекомендации и перспективы дальнейшей разработки темы.
%\begin{enumerate}
  %\item На основе анализа \ldots
  %\item Численные исследования показали, что \ldots
  %\item Математическое моделирование показало \ldots
  %\item Для выполнения поставленных задач был создан \ldots
%\end{enumerate}

Представленная диссертация предлагает обзор морфоструктур 20 желобов Тихого океана, выполненный с использованием передовых методов анализа данных, моделирования и картографирования. Исследование представляет собой междисциплинарный подход, сочетающий картографическое моделирование и тектонико-геологический региональный анализ. Результаты позволили представить вариативность желобов на основе их геоморфологических свойств, изменения уклонов, анализ градиентов крутизны и характер батиметрии. Анализ данных в геотектонике и геологии с использованием статистических и скриптовых картографических методов является важным подходом, направленным на повышение точности обработки больших массивов данных. Методы, представленные в диссертации, направлены на решение проблемы анализа датасетов для вопросов моделирования в морской геоморфологии, что связано с тем, что существующие решения в основном применяют традиционные методы ГИС. Диссертация продемонстрировала использование методов скриптовой картографии применительно к задачам геотектоники и геологии для  анализа подводных структур рельефа дна на примере 20 желобов Тихого океана. 
